\section{曼哈顿最小生成树：闭锥扫描与平局割证明}\label{knowledge-manhattan-mst}
\index{Manhattan MST}\index{minimum spanning tree}
接口见第~\ref{compact-ManhattanMST}~节（第~\pageref{compact-ManhattanMST}~页），
正式题用法见第~\ref{usage-example-236}~节（第~\pageref{usage-example-236}~页）。
完整无向图的边权是 $d(p,q)=|p_x-q_x|+|p_y-q_y|$。
不枚举二次边集：每个点只保留四个向右闭锥中的最近邻，再对候选边作 Kruskal。
候选图含有至少一棵最小生成树，不保证含所有最优边；禁边、颜色等额外约束不能直接套用。

\paragraph{一个方向的扫描}
变换后的锥取 $q_x-p_x\ge q_y-p_y\ge0$。
按 $(x+y,\text{原编号})$ 排序，活动点用 $-y$ 作 map 键。
依次从 $y_p\le y_q$ 的点开始检查，直到 $x_q-x_p<y_q-y_p$。
活动点的 $x-y$ 随 $y$ 递减而严格递增，因此首次失败后可停止。
第一次满足条件的后继具有最小 $x+y$，其距离正是 $(x_q+y_q)-(x_p+y_p)$。
删掉已找到邻居的活动点，每点每轮最多删除一次，故每轮至多 $n-1$ 条候选边。
四轮坐标分别为 $(x,y),(y,x),(-y,x),(x,-y)$，覆盖右侧四个闭锥。
任意无向边都能从左端点指向其中一个锥；竖直边任选一个方向。
总时间 $O(n\log n)$，额外空间 $O(n)$。

\paragraph{重复坐标与闭边界}
重复坐标是不同顶点。排序平局按原编号，连续代表之间生成零权链；
同一 $x+y$ 的其他位置不能以非零距离删掉该代表，最后一个代表继续寻找非重合邻居。
可先在证明中收缩这些零链。若一个割拆开重合点类，零链已有最小的过割边；否则割可直接下降到收缩图。
闭锥内只能使用 $d(r,q)\le\max(d(p,r),d(p,q))$，不能写严格小于：
$p=(0,0),q=(2,0),r=(1,1)$ 的三条距离都是2。

\paragraph{平局也成立的割证明}
在非重合点图的任意割中，选距离最小为 $\delta$ 的过割点对 $p,q$，
并在平局中令 $p_x+q_x$ 最大；令 $p$ 为左端点。
选择包含 $q$ 的右向锥，若边恰在对角线上，选靠近水平轴的浅锥。
该锥的最近邻 $r$ 满足 $d(p,r)\le\delta$。
若 $r$ 和 $q$ 在割的同侧，则 $pr$ 已是候选最小过割边。
否则 $rq$ 过割且 $d(r,q)\le\delta$：严格小于违反最小性，等于也会违反上述平局选择。

后一个结论可在局部锥坐标 $X\ge Y\ge0$ 中检查。
满足 $d(p,r)\le\delta$ 的点位于顶点
$(0,0),(\delta,0),(\delta/2,\delta/2)$ 的三角形。
当 $q$ 严格位于锥内时，距离达到 $\delta$ 只能取已排除的原点；
$q$ 在主轴时，非零等号点在 $X=Y$ 的射线上；
$q$ 在对角线时，非零等号点只能是 $(\delta,0)$。
后两种情况在所选右向锥中都使 $r_x>p_x$（对角线选浅锥正是为了保证这一点），
于是过割点对 $r,q$ 的横坐标和更大，矛盾。
所以每个割都保留最小边；在候选图上应用 Prim 的割性质，得到的树也是完整图的 MST。

\paragraph{整数范围与抄写合同}
输入为 signed 64 位整数对；先整体转换为 signed 128 位，再进行和、差、变号及距离求和。
单边最大距离为 $2^{65}-2$；在通常32位 int 编号上限内，总权远小于 signed 128 位上限。
返回原始0基编号和128位总权；空集与单点返回0和空边集。
正式题限制 $1\le n\le200000$、$0\le x,y\le10^9$，仅该驱动可安全转回 long long 输出。
并查集只在本调用点按大小选择根的合并顺序，不改普通 dsu 的接口语义。

参考：Zhou--Shenoy--Nicholls (2002)，DOI \texttt{10.1016/S0020-0190(01)00232-0}；
四轮 map 实现参考 KACTL \texttt{ManhattanMST.h}（CC0，固定版本与字节哈希见来源记录）。
原论文半开锥的严格不等式不能直接代入这里的闭锥代码。
